\documentclass{article}
\usepackage{graphicx} % Required for inserting images
\usepackage{amsmath}
\usepackage{algorithmic}
\usepackage{algorithm}
\usepackage[polish]{babel}
\usepackage[utf8]{inputenc}
\usepackage[T1]{fontenc}
\usepackage{float}

\title{Algorytmy Równoległe - Laboratorium 1}
\author{Adam Staniszewski}

\begin{document}

\maketitle

\section*{Wybór problemu}
Wyznaczyć temperaturę w przewodniku elektrycznym, który został nagrzany w wyniku przepływu prądu elektrycznego o stałym natężeniu (powierzchnia boczna jest utrzymywana w stałej temperaturze równej zero).
Zastosować metodę różnicową do równania Poissona:

\[
\frac{\partial^2 T}{\partial x^2} + \frac{\partial^2 T}{\partial y^2} = -\frac{g}{\lambda}
\]

gdzie:
- \( g \) - wydajność źródeł ciepła (stała),
- \( \lambda \) - przewodność cieplna (stała).


\section*{Wyprowadzenie wzoru iteracyjnego}
Warunki brzegowe - powierzchnia boczna jest utrzymywana w stałej temperaturze równej zero.
Równanie Poissona dla rozkładu temperatury zapisujemy jako:
\[
\frac{\partial^2 T}{\partial x^2} + \frac{\partial^2 T}{\partial y^2} = -\frac{g(x, y)}{\lambda}
\]
gdzie:
\begin{itemize}
    \item \( T(x, y) \) - funkcja temperatury,
    \item \( g(x, y) \) - funkcja źródła ciepła,
    \item \( \lambda \) - przewodność cieplna.
\end{itemize}
Przybliżamy drugie pochodne różnicami skończonymi. Dla drugiej pochodnej względem \( x \) mamy:
\[
\frac{\partial^2 T}{\partial x^2} \approx \frac{T_{i+1, j} - 2T_{i, j} + T_{i-1, j}}{dx^2}
\]
Analogicznie, dla drugiej pochodnej względem \( y \):
\[
\frac{\partial^2 T}{\partial y^2} \approx \frac{T_{i, j+1} - 2T_{i, j} + T_{i, j-1}}{dy^2}
\]
Sumujemy teraz oba przybliżenia, co daje:
\[
\frac{T_{i+1, j} - 2T_{i, j} + T_{i-1, j}}{dx^2} + \frac{T_{i, j+1} - 2T_{i, j} + T_{i, j-1}}{dy^2} = -\frac{g_{i, j}}{\lambda}
\]
Zakładając, że siatka jest jednolita, czyli \( dx = dy \), możemy uprościć to wyrażenie do postaci:
\[
\frac{T_{i+1, j} + T_{i-1, j} + T_{i, j+1} + T_{i, j-1} - 4T_{i, j}}{dx^2} = -\frac{g_{i, j}}{\lambda}
\]
Rozwiązując to równanie względem \( T_{i, j} \), otrzymujemy wzór na nową wartość temperatury:
\[
T_{\text{new}}[i][j] = \frac{1}{4} \left( T_{i+1, j} + T_{i-1, j} + T_{i, j+1} + T_{i, j-1} + \frac{g_{i, j} \cdot dx^2}{\lambda} \right)
\]

\section*{Pseudokod algorytmu sekwencyjnego}
Algorytm rozwiązuje równanie Poissona metodą różnic skończonych dla siatki 2D, aby wyznaczyć rozkład temperatury \(T\) w przewodniku elektrycznym. 

\begin{algorithm}[H]
\caption{Algorytm sekwencyjny}
\begin{algorithmic}[1]
    \STATE Inicjalizacja:
    \STATE Definicja parametrów:
    \STATE \quad $g$ - wydajność źródła ciepła (stała)
    \STATE \quad $\lambda$ - przewodność cieplna (stała)
    \STATE \quad $dx, dy$ - odległość między węzłami siatki w kierunkach $x$ i $y$
    \STATE \quad $\epsilon$ - tolerancja dla kryterium zbieżności
    \STATE \quad max\_iterations - maksymalna liczba iteracji
    \STATE \quad $T[x][y]$ - macierz temperatury
    \STATE Zainicjuj macierz $T[x][y]$ zerami (warunki brzegowe: temperatura = 0)
    \REPEAT
    \FOR {każdy punkt wewnętrzny $(i, j)$}
        \STATE Oblicz nową wartość temperatury:
        \[
        T_{\text{new}}[i][j] = \frac{1}{4} (T[i+1][j] + T[i-1][j] + T[i][j+1] + T[i][j-1] + \frac{g \times dx^2}{\lambda})
        \]
    \ENDFOR
    \STATE Sprawdź zbieżność:
    \IF {różnica $|T_{\text{new}}[i][j] - T[i][j]|$ dla każdego $(i,j)$ jest mniejsza niż $\epsilon$}
        \STATE Zbieżność osiągnięta, zakończ iterację
    \ENDIF
    \STATE Zaktualizuj macierz $T$: $T = T_{\text{new}}$
    \UNTIL{osiągnięcie zbieżności lub maksymalnej liczby iteracji}
    \STATE Wyświetl wynikową macierz $T$
\end{algorithmic}
\end{algorithm}

\section*{Opis poszczególnych faz PCAM}
\begin{itemize}
    \item \textbf{P}artition - w jaki sposób rozdzielić zasoby między procesorami? W tym przypadku możemy podzielić siatkę na mniejsze cześci obsługiwane przez pojedyncze procesory.
    \item \textbf{C}ommunicate - w jaki sposób procesory wymieniają informacje? Procesory w tym przypadku mogą wymieniać dane brzegowe i wyniki obliczeń.
    \item \textbf{A}gglomerate - analiza zadań oraz struktury z poprzednich punktów. W razie potrzeby można łączyć zadania w większe, aby zmniejszyć narzut komunikacyjny.
    \item \textbf{M}apping - przypisanie zadań do poszczególnych procesorów. W naszym przypadku możemy na przykład ustalić, że procesor 0 będzie zbierał wyniki z innych procesorów i sprawdzał czy została osiągnięta zbieżność.
\end{itemize}

\section*{Psuedokod algorytmu równoległego z wykorzystaniem metodyki PCAM}
Algorytm równoległy dzieli siatkę na bloki przetwarzane równolegle przez różne procesory z wymianą danych granicznych.

\begin{algorithm}[H]
\caption{Algorytm równoległy (metoda różnic skończonych z MPI)}
\begin{algorithmic}[1]
    \STATE Podziel domenę na $n$ części, gdzie $n$ to liczba procesów.
    \STATE Każdy proces otrzymuje lokalny blok danych, zawierający:
    \STATE \quad Obszar obliczeniowy $T_{\text{local}}[x][y]$ oraz dane potrzebne do obliczeń: $g$, $\lambda$, $dx$, $dy$.
    
    \STATE Uruchom MPI, zdefiniuj $rank$ (numer procesu) oraz $size$ (liczba procesów)
    \STATE Każdy proces otrzymuje swój blok danych, w tym:
    \STATE \quad $g$, $\lambda$, $dx$, $dy$, $T_{\text{local}}[x][y]$
    
    \REPEAT
    \FOR {każdy punkt wewnętrzny $(i, j)$ w lokalnym bloku}
        \STATE Oblicz nową wartość temperatury:
        \[
        T_{\text{new}}[i][j] = \frac{1}{4} (T_{\text{local}}[i+1][j] + T_{\text{local}}[i-1][j] + T_{\text{local}}[i][j+1] + T_{\text{local}}[i][j-1] + \frac{g \times dx^2}{\lambda})
        \]
    \ENDFOR
    
    \STATE Wymień dane graniczne z sąsiadującymi procesami (zbiorczo, w jednym przesłaniu)
    
    \STATE Każdy proces oblicza lokalną różnicę $|T_{\text{new}}[i][j] - T_{\text{local}}[i][j]|$

    \STATE Zbierz różnice na procesie 0 (root) za pomocą $MPI\_Reduce$
    
    \IF {globalna różnica jest mniejsza niż $\epsilon$ (na procesie 0)}
    \STATE Zbieżność osiągnięta, zakończ iterację
    \ENDIF
    \STATE Zaktualizuj macierz $T$: $T_{\text{local}} = T_{\text{new}}$
    \UNTIL{osiągnięcie zbieżności lub maksymalnej liczby iteracji}
    
    \STATE Proces 0 zbiera wyniki z wszystkich procesów za pomocą $MPI\_Gather$

    \STATE Wyświetl wynikową macierz $T$ (na procesie 0)
    
\end{algorithmic}
\end{algorithm}

\end{document}
